\section{Proofs of the local law and of the first-order bounds}\label{app:torus}

This appendix proves Proposition~\ref{prop:constant}(a),
Lemma~\ref{lem:upper}, Theorem~\ref{thm:local}, Lemma~\ref{lem:discrete} and
Theorem~\ref{thm:first}(ii) for reducible faces. Except in
Section~\ref{app:first}, $F$ is a face with $k\ge2$ states, numbered
$1,\dots,k$. For $x\in\mathbb C^F$ we write $x'=(x_2,\dots,x_k)$. We write
$\mathrm i$ for the imaginary unit, since $i$ denotes a state, and
$\mathbb T^m=(\R/2\pi\Z)^m$. We use 2 standard facts~\cite{bermanplemmons1994}.
First, let $B\ge0$ be square and $S$ diagonal with positive diagonal. Then
$\varrho(B-S)<0$ if and only if $\varrho(S^{-1}B)<1$, and also if and only if
$(S-B)v>0$ for some $v>0$. In that case
$(S-B)^{-1}=\sum_{j\ge0}(S^{-1}B)^jS^{-1}\ge0$. Second, if $M$ is an
irreducible Metzler matrix, then $\varrho(M)$ is larger than the Perron root
of each proper principal submatrix of $M$, and every nonnegative eigenvector
of $M$ is positive and belongs to $\varrho(M)$.

\subsection{The Perron root under a diagonal perturbation}

An irreducible generator on $F$ is a Metzler matrix $G$ with $G\ones=0$ whose
graph is strongly connected. It has a unique stationary law $\pi>0$. The
argument after Definition~\ref{def:tilt}, with $(G,\pi)$ in place of
$(\widehat Q_\alpha,\alpha)$, shows that
$Z=(\ones\pi^\top-G)^{-1}-\ones\pi^\top$ exists and
satisfies~\eqref{eq:Zprops}. Put $D=\operatorname{diag}\pi$ and
$\Gamma=DZ+Z^\top D$.

\begin{lemma}[the Perron root]\label{lem:perron}
Let $G$, $\pi$, $Z$ and $\Gamma$ be as above. Then
$\kappa(\eta)=\varrho(G+\operatorname{diag}\eta)$ is real-analytic near $0$,
with $\kappa(0)=0$, $\nabla\kappa(0)=\pi$ and $\nabla^2\kappa(0)=\Gamma$. For
$\eta\in\R^F$ and $u=Z\eta$,
\begin{equation}\label{eq:A-dirichlet}
\eta^\top\Gamma\eta=\sum_{i\ne j}\pi_iG_{ij}(u_i-u_j)^2 .
\end{equation}
So $\Gamma$ is positive semidefinite with kernel $\R\ones$, and $\Sigma$,
which is $\Gamma$ without the row and column of state $1$, is positive
definite.
\end{lemma}

\begin{proof}
Since $G+\operatorname{diag}\eta$ is irreducible, its Perron root is a simple
eigenvalue. So by the implicit function theorem, $\kappa$ and the right
eigenvector $v(\eta)$ with $\pi^\top v=1$ are real-analytic near $0$, with
$\kappa(0)=0$ and $v(0)=\ones$. With subscripts for derivatives at $\eta=0$,
differentiating
$(G+\operatorname{diag}\eta)v=\kappa v$ in $\eta_j$ gives
$Gv_j+e_j=\kappa_j\ones$, and multiplying by $\pi^\top$ gives
$\kappa_j=\pi_j$. Then $v_j=Ze_j$, because by~\eqref{eq:Zprops} both vectors
solve $Gx=\pi_j\ones-e_j$ with $\pi^\top x=0$, and the kernel of $G$ is
$\R\ones$. Differentiating again, in $\eta_i$, gives
$Gv_{ij}+\operatorname{diag}(e_i)v_j+\operatorname{diag}(e_j)v_i
=\kappa_{ij}\ones+\kappa_iv_j+\kappa_jv_i$, and multiplying by $\pi^\top$
gives $\kappa_{ij}=\pi_iZ_{ij}+\pi_jZ_{ji}=\Gamma_{ij}$.

For~\eqref{eq:A-dirichlet} let $f=\eta-(\pi\cdot\eta)\ones$.
By~\eqref{eq:Zprops}, $-Gu=f$ and $\pi\cdot u=0$, so
$\eta^\top\Gamma\eta=2\sum_i\pi_i\eta_iu_i=2\sum_i\pi_if_iu_i
=-2\sum_i\pi_iu_i(Gu)_i$. Expanding the squares in~\eqref{eq:A-dirichlet}
and using $G\ones=0$ and $\pi^\top G=0$ gives the same expression. Also
$\Gamma\ones=0$ by~\eqref{eq:Zprops}. If $\eta^\top\Gamma\eta=0$, then $u$ is
constant along every edge of the strongly connected graph of $G$, so $u$ is
constant. Then $u=0$ because $\pi\cdot u=0$, so $f=0$ and $\eta\in\R\ones$. A
nonzero vector that vanishes at state $1$ is not in $\R\ones$, which gives
the claim about $\Sigma$.
\end{proof}

\begin{proof}[Proof of Proposition~\ref{prop:constant}(a)]
For $\vartheta\in\R^F$ let
$\Lambda(\vartheta)=\varrho(Q_F+\operatorname{diag}\vartheta)$, so that
$\Lambda(\vartheta+c\ones)=\Lambda(\vartheta)+c$. It is a simple eigenvalue
with positive right and left eigenvectors $\hat r,\hat l$, normalized by
$\sum_i\hat r_i=1$ and $\hat l\cdot\hat r=1$, and by the implicit function
theorem $\Lambda$, $\hat r$ and $\hat l$ are real-analytic in $\vartheta$. Put
$\pi=\hat l\circ\hat r$ and
$G=\hat R^{-1}(Q_F+\operatorname{diag}\vartheta-\Lambda(\vartheta)I)\hat R$
with $\hat R=\operatorname{diag}\hat r$. Then $G$ is an irreducible generator
with stationary law $\pi$. Since $\hat R$ commutes with diagonal matrices,
$\Lambda(\vartheta+\eta)=\Lambda(\vartheta)+\varrho(G+\operatorname{diag}\eta)$.
So by Lemma~\ref{lem:perron}, $\nabla\Lambda(\vartheta)=\pi$ and
$\nabla^2\Lambda(\vartheta)=\Gamma$, with $\Gamma$ built from $(G,\pi)$.

Let $\Lambda_0(\beta)=\Lambda(0,\beta)$ for $\beta\in\R^{k-1}$. Its Hessian
is $\Gamma$ without the row and column of state $1$, which is positive
definite by Lemma~\ref{lem:perron}. So $\nabla\Lambda_0$ is injective, and by
the inverse function theorem it has a real-analytic inverse on its open
image. Fix $\alpha\in\Delta_F^\circ$ and put
$\vartheta(\alpha)=q-\theta(\alpha)$. Then
$Q_F+\operatorname{diag}\vartheta(\alpha)=A_F-\operatorname{diag}\theta(\alpha)$
has Perron root $0$ and, by~\eqref{eq:stationary}, Perron vectors $r_\alpha$
and $l_\alpha$. So at $\vartheta(\alpha)$ we have $\pi=\alpha$,
$G=\widehat Q_\alpha$ and $\Gamma=\Gamma_\alpha$. Let $\beta(\alpha)$ be the
last $k-1$ coordinates of $\vartheta(\alpha)-\vartheta_1(\alpha)\ones$. By the
shift property, $\nabla\Lambda_0(\beta(\alpha))=\alpha'$ and
$\nabla^2\Lambda_0(\beta(\alpha))=\Sigma_\alpha$. Hence
$\beta(\alpha)=(\nabla\Lambda_0)^{-1}(\alpha')$ is real-analytic, and so are
$\vartheta_1(\alpha)=-\Lambda_0(\beta(\alpha))$, $\theta(\alpha)$, its
Perron vectors and $\Sigma_\alpha$. Moreover
$I_F(\alpha)=\langle\vartheta(\alpha),\alpha\rangle
=\langle\beta(\alpha),\alpha'\rangle-\Lambda_0(\beta(\alpha))$.
Differentiating and using $\nabla\Lambda_0(\beta)=\alpha'$ gives
$\nabla I_F=\beta$. So $\nabla^2I_F$ is the Jacobian of
$(\nabla\Lambda_0)^{-1}$ at $\alpha'$, which is $\Sigma_\alpha^{-1}$. The
formula for $C_F$ follows from Definition~\ref{def:tilt}.
\end{proof}

\subsection{The torus representation}

We split
\[
A_F=\begin{pmatrix}0&a_1^\top\\ b_1&A'\end{pmatrix},\qquad
\mu_F=(\mu_1,\mu'),
\]
so $a_1=(q_{1j})_{j\ge2}$, $b_1=(q_{i1})_{i\ge2}$ and
$A'=(q_{ij})_{i,j\ge2}$. For $z'\in\mathbb C^{k-1}$ let
$\mathcal Z'=\operatorname{diag}(z')$. Where $\mathcal Z'-A'$ is invertible
we put $X(z')=(\mathcal Z'-A')^{-1}$, $g(z')=a_1^\top X(z')b_1$ and
\[
b(z')=\bigl(\mu_1+\mu'^\top X(z')b_1\bigr)\bigl(1+a_1^\top X(z')\ones\bigr).
\]

\begin{lemma}[torus representation]\label{lem:torus}
Let $s'\in(0,\infty)^{k-1}$ satisfy
$\varrho(A'-\operatorname{diag}s')<0$. On the set of $z'\in\mathbb C^{k-1}$
with $|z_i|\ge s_i$ for every $i\ge2$, the matrix
$\mathcal Z'-A'$ is invertible, and $X$, $g$ and $b$ are given by absolutely
convergent power series in $1/z_2,\dots,1/z_k$ with nonnegative coefficients.
In particular $|g(z')|\le g(s')$ and
$|b(z')|\le b(s')$ there. For every $t\in(0,\infty)^F$,
\[
H(t)=\frac1{(2\pi)^{k-1}}\int_{\mathbb T^{k-1}}
\exp\Bigl(t_1g(z')+\sum_{i\ge2}t_iz_i\Bigr)\,b(z')\prod_{i\ge2}z_i\,d\varphi,
\qquad z_i=s_ie^{\mathrm i\varphi_i}.
\]
\end{lemma}

\begin{proof}
By the first fact, $\varrho(\operatorname{diag}(s')^{-1}A')<1$. So the series
$X(z')=\sum_{j\ge0}(\mathcal Z'^{-1}A')^j\mathcal Z'^{-1}$ converges
absolutely for $|z_i|\ge s_i$, since its terms are bounded
entrywise in modulus by those at $z'=s'$. As $a_1$, $b_1$ and $\mu$ are
nonnegative, this gives the claims about $X$, $g$ and $b$.

Choose $s_1>g(s')$ and put $S=\operatorname{diag}(s_1,s')$. We show that
$\varrho(S^{-1}A_F)<1$. For $\varepsilon>0$ let $v=(1,v')$ with
$v'=X(s')(b_1+\varepsilon\ones)>0$. Then $((S-A_F)v)_i=\varepsilon$ for
$i\ge2$, and $((S-A_F)v)_1=s_1-g(s')-\varepsilon a_1^\top X(s')\ones$. For
small $\varepsilon$ both are positive, and the first fact gives the claim.

Now let $\mathcal Z=\operatorname{diag}(z)$ with $|z_i|=s_i$ for
every $i\in F$. Since $\varrho(A_FS^{-1})<1$, expanding
$(\mathcal Z-A_F)^{-1}=\mathcal Z^{-1}(I-A_F\mathcal Z^{-1})^{-1}$ gives
\[
\mu_F^\top(\mathcal Z-A_F)^{-1}\ones=\sum_c\mu_{c_1}\prod_{j<M}q_{c_jc_{j+1}}
\prod_{i\in F}z_i^{-m_i(c)},
\]
absolutely and uniformly on the torus. Here $c=c_1\cdots c_M$ runs over the
words in $F$ with $M\ge1$ and $c_j\ne c_{j+1}$, and $m_i(c)$ is the number of
letters $i$ in $c$. For an integer $m$ and $t_i>0$, Cauchy's formula gives
$\frac1{2\pi}\int_0^{2\pi}e^{t_iz_i}z_i^{1-m}\,d\varphi_i=t_i^{m-1}/(m-1)!$ if
$m\ge1$ and $0$ if $m\le0$. So integrating term by term removes the words
that miss a state of $F$, and grouping the others by $m(c)$ gives
\begin{equation}\label{eq:A-full}
H(t)=\frac1{(2\pi)^k}\int_{\mathbb T^k}e^{\langle t,z\rangle}\,
\mu_F^\top(\mathcal Z-A_F)^{-1}\ones\prod_{i\in F}z_i\,d\varphi .
\end{equation}
The Schur complement $z_1-g(z')$ of $\mathcal Z'-A'$ in $\mathcal Z-A_F$ is
not $0$, since $|g(z')|\le g(s')<s_1$. Block inversion gives
\[
\mu_F^\top(\mathcal Z-A_F)^{-1}\ones=\frac{b(z')}{z_1-g(z')}
+\mu'^\top X(z')\ones .
\]
The second term does not depend on $z_1$. So by Cauchy's formula the
integral over $\varphi_1$ in~\eqref{eq:A-full}, divided by $2\pi$, is
$e^{t_1g(z')}b(z')$ times the factors without $z_1$. This proves the
formula.
\end{proof}

All coefficients are nonnegative, so the integrand has its largest modulus
at $\varphi=0$. This gives Lemma~\ref{lem:upper}, and a Laplace expansion
around $\varphi=0$ gives Theorem~\ref{thm:local}.

\subsection{Proof of \texorpdfstring{Lemma~\ref{lem:upper}}{the upper bound}}

Let $F$ be irreducible and $\theta\in V_F$. Then
$A_F-\operatorname{diag}\theta$ is an irreducible Metzler matrix with Perron
root $0$. So it has a right eigenvector $u>0$ for $0$, and
$\theta_i=(A_Fu)_i/u_i$. Each row of $A_F$ has a positive entry, since $G_F$
is strongly connected and $k\ge2$, so $\theta>0$. By the second fact,
$\varrho(A'-\operatorname{diag}\theta')<0$. The rows $i\ge2$ of
$(A_F-\operatorname{diag}\theta)u=0$ give $u'=u_1X(\theta')b_1$, and then row
$1$ gives $\theta_1=g(\theta')$. Hence
\begin{equation}\label{eq:A-radii}
\theta>0,\qquad \varrho(A'-\operatorname{diag}\theta')<0,\qquad
g(\theta')=\theta_1\qquad\text{for every }\theta\in V_F .
\end{equation}

Now fix $n$ with support $F$ and $h$ with $hq_{\max}\le\frac12$, and put
$|n|=\sum_in_i$, $t=hn$ and $\sigma=e^{hq_{\max}}$.
Lemma~\ref{lem:runs} with $|n|$ in place of $N$ gives the
probability $P_{|n|}(n)$ in the lemma. Only terms with $m_i\le n_i$
for every $i$ are nonzero, and in them
$h^{m_i-1}\binom{n_i-1}{m_i-1}\le t_i^{m_i-1}/(m_i-1)!$ and
$(1-hq_i)^{n_i-m_i}\le e^{-q_it_i}\sigma^{m_i}$. Note that
$h^{M-1}=h^{k-1}\prod_ih^{m_i-1}$ and $\sigma^MW(m)=\sigma W_\sigma(m)$, where
$W_\sigma$ is $W$ for $\sigma A_F$ in place of $A_F$. So
\begin{equation}\label{eq:A-sigma}
P_{|n|}(n)\le\sigma\,h^{k-1}e^{-\langle q,t\rangle}H_\sigma(t),
\end{equation}
where $H_\sigma$ is $H$ for $\sigma A_F$. We bound $H_\sigma$ by
Lemma~\ref{lem:torus} for $\sigma A_F$ with the radii $\sigma\theta'$, which is
allowed since
$\varrho(\sigma A'-\operatorname{diag}\sigma\theta')
=\sigma\varrho(A'-\operatorname{diag}\theta')<0$. The functions of
Lemma~\ref{lem:torus} for $\sigma A_F$ satisfy
$X_\sigma(\sigma z')=\sigma^{-1}X(z')$, so $g_\sigma(\sigma z')=\sigma g(z')$
and $b_\sigma(\sigma z')=b(z')$. By~\eqref{eq:A-radii} the integrand has
modulus at most $e^{\sigma\langle t,\theta\rangle}b(\theta')\sigma^{k-1}
\prod_{i\ge2}\theta_i$, so for every $\sigma>0$
\begin{equation}\label{eq:A-Hsigma}
H_\sigma(t)\le\sigma^{k-1}b(\theta')\prod_{i\ge2}\theta_i\,
e^{\sigma\langle t,\theta\rangle}.
\end{equation}
With~\eqref{eq:A-sigma} this gives
$P_{|n|}(n)\le\sigma^kb(\theta')\prod_{i\ge2}\theta_i\,h^{k-1}
\exp(-h\langle n,q-\theta\rangle+(\sigma-1)h\langle n,\theta\rangle)$.
Finally $\sigma\le e^{1/2}$ and $\sigma-1\le2hq_{\max}$, so the lemma holds
with $C_\theta=\max\bigl(e^{k/2}b(\theta')\prod_{i\ge2}\theta_i,\
2q_{\max}\max_i\theta_i\bigr)$. The bound after the lemma is the case
$\theta=q-\lambda_F\ones$. \qed

\subsection{Proof of \texorpdfstring{Theorem~\ref{thm:local}}{the local law}}

For $k=1$ the formula is part of Lemma~\ref{lem:continuum}. Let $k\ge2$ and
$\alpha\in\mathcal K$, and write $\theta=\theta(\alpha)$, $r=r_\alpha$,
$l=l_\alpha$ and $\Theta'=\operatorname{diag}(\theta')$. By~\eqref{eq:A-radii}
we may use Lemma~\ref{lem:torus} with the radii $\theta'$ and $t=T\alpha$.
With Lemma~\ref{lem:continuum} this gives
\begin{equation}\label{eq:A-fT}
f_T(\alpha)=\frac{T^{k-1}e^{-T\langle q,\alpha\rangle}}{(2\pi)^{k-1}}
\int_{\mathbb T^{k-1}}e^{T\Psi(\alpha,\varphi)}\chi(\alpha,\varphi)\,d\varphi,
\end{equation}
where $z_i=\theta_ie^{\mathrm i\varphi_i}$,
$\Psi(\alpha,\varphi)=\alpha_1g(z')+\sum_{i\ge2}\alpha_iz_i$ and
$\chi(\alpha,\varphi)=b(z')\prod_{i\ge2}z_i$. We check the hypotheses of
Lemma~\ref{lem:laplace} with $m=k-1$.

\emph{Hypothesis \textup{(iii)}.} By~\eqref{eq:A-radii},
$\Psi(\alpha,0)=\langle\alpha,\theta\rangle$ is real. By
Lemma~\ref{lem:torus}, $|g(z')|\le g(\theta')=\theta_1$ on the
torus, so
\[
\operatorname{Re}\Psi(\alpha,\varphi)\le\Psi(\alpha,0)
-\sum_{i\ge2}\alpha_i\theta_i(1-\cos\varphi_i).
\]
All $\alpha_i$ and $\theta_i$ are positive, so the sum is positive unless
$\varphi=0$.

\emph{Hypotheses \textup{(i)} and \textup{(ii)}.} For $z\in\R^F$ near
$\theta$ let $\lambda(z)=\varrho(A_F-\operatorname{diag}z)$. With
$R=\operatorname{diag}r$ we have
$A_F-\operatorname{diag}(\theta-\eta)=R(\widehat
Q_\alpha+\operatorname{diag}\eta)R^{-1}$,
and $\widehat Q_\alpha$ is an irreducible generator with stationary law
$\alpha$. So Lemma~\ref{lem:perron} gives $\nabla\lambda(\theta)=-\alpha$ and
$\nabla^2\lambda(\theta)=\Gamma_\alpha$. For real $z'$ near $\theta'$ we have
$\varrho(A'-\operatorname{diag}z')<0$, so $(1,X(z')b_1)\ge0$ by the first
fact. By the definition of $g$ this vector lies in the kernel of
$A_F-\operatorname{diag}(g(z'),z')$, so $\lambda(g(z'),z')=0$ by the second
fact. Write $g_j$, $g_{ij}$ for the derivatives of $g$ at $\theta'$ and put
$w_j=e_j-(\alpha_j/\alpha_1)e_1\in\R^F$ for $j\ge2$. Differentiating
$\lambda(g(z'),z')=0$ once and twice at $\theta'$ gives
\[
g_j=-\frac{\alpha_j}{\alpha_1},\qquad \alpha_1g_{ij}=w_i^\top\Gamma_\alpha w_j .
\]
So $\partial\Psi/\partial z_j=\alpha_1g_j+\alpha_j=0$ at $\theta'$, and the
Hessian $\mathcal H$ of $\Psi$ in $z'$ at $\theta'$ is $V^\top\Gamma_\alpha V$,
where $V$ has columns $w_2,\dots,w_k$. Since
$\partial/\partial\varphi_j=\mathrm iz_j\,\partial/\partial z_j$ and the first
derivatives vanish, $\nabla_\varphi\Psi(\alpha,0)=0$ and
$-\nabla^2_\varphi\Psi(\alpha,0)=\Theta'\mathcal H\Theta'$. Let
$B=I+\ones\alpha'^\top/\alpha_1$. Then
$Vy+(\alpha'^\top y/\alpha_1)\ones=(0,By)$ for $y\in\R^{k-1}$, and
$\Gamma_\alpha\ones=0$, so $\mathcal H=B^\top\Sigma_\alpha B$. By
Lemma~\ref{lem:perron}, $\Sigma_\alpha$ is positive definite, and hence so is
$\mathcal H$. The matrix determinant lemma gives $\det B=1/\alpha_1$, so
$\det\mathcal H=\det\Sigma_\alpha/\alpha_1^2$.

\emph{Regularity.} By Proposition~\ref{prop:constant}(a), $\theta(\alpha)$ is
real-analytic. Fix $\alpha_0\in\mathcal K$. By~\eqref{eq:A-radii} and
continuity there is $\delta\in(0,\min_i\theta_i(\alpha_0))$ with
$\varrho(A'-\operatorname{diag}(\theta'(\alpha_0)-\delta\ones))<0$. By
Lemma~\ref{lem:torus}, $X$, $g$ and $b$ are holomorphic where
$|z_i|>\theta_i(\alpha_0)-\delta$ for every $i\ge2$, and this holds
on the torus of radii $\theta'(\alpha)$ for $\alpha$ near $\alpha_0$. So $\Psi$
and $\chi$ are real-analytic near $\{\alpha_0\}\times\mathbb T^{k-1}$.
Covering $\mathcal K$ by finitely many such sets shows that their
$\varphi$-derivatives are jointly continuous on
$\mathcal K\times\mathbb T^{k-1}$.

\emph{The constant.} Lemma~\ref{lem:laplace} now gives, uniformly for
$\alpha\in\mathcal K$ and $T\ge1$,
\[
\int_{\mathbb T^{k-1}}e^{T\Psi}\chi\,d\varphi
=e^{T\langle\alpha,\theta\rangle}\Bigl(\frac{2\pi}T\Bigr)^{\frac{k-1}2}
\frac{\alpha_1}{\prod_{i\ge2}\theta_i\sqrt{\det\Sigma_\alpha}}
\Bigl(\chi(\alpha,0)+O\Bigl(\frac1T\Bigr)\Bigr).
\]
Here $\chi(\alpha,0)=b(\theta')\prod_{i\ge2}\theta_i$. As in the derivation
of~\eqref{eq:A-radii}, $r'=r_1X(\theta')b_1$. In the same way the columns
$j\ge2$ of $l^\top(A_F-\operatorname{diag}\theta)=0$, which is
\eqref{eq:stationary}, give $l'=l_1X(\theta')^\top a_1$. Hence
$\mu\cdot r=r_1(\mu_1+\mu'^\top X(\theta')b_1)$ and
$l\cdot\ones=l_1(1+a_1^\top X(\theta')\ones)$. Since $l_1r_1=\alpha_1$,
\[
b(\theta')=\frac{(\mu\cdot r_\alpha)(l_\alpha\cdot\ones)}{\alpha_1}.
\]
As $r_\alpha,l_\alpha>0$ and $\mu(F)>0$ by (H), $\chi(\alpha,0)$ is positive
and continuous, hence bounded below on $\mathcal K$. So the last bracket is
$\chi(\alpha,0)(1+O(1/T))$. Putting this into~\eqref{eq:A-fT} and using
$I_F(\alpha)=\langle q-\theta,\alpha\rangle$ gives
$f_T(\alpha)=T^{(k-1)/2}C_F(\alpha)e^{-TI_F(\alpha)}(1+O(1/T))$, uniformly on
$\mathcal K$. \qed

\subsection{Proof of \texorpdfstring{Lemma~\ref{lem:discrete}}{the discrete
law}}

Under the continuum weights a path has of order $T$ runs, and each run after
the first in state $i$ changes the number of placements by a factor
$1-O(m_i/n_i)$, which suggests a relative error of order $T^2/N$. Let $t=hn$
and $\alpha=n/N\in\mathcal K$. By Lemma~\ref{lem:continuum},
$N^{-(k-1)}f_T(n/N)=h^{k-1}e^{-\langle q,t\rangle}H(t)=\sum_mv_m$ with
\[
v_m=W(m)\,h^{k-1}\prod_{i\in F}e^{-q_it_i}\frac{t_i^{m_i-1}}{(m_i-1)!},
\qquad m\in\Z_{\ge1}^F .
\]
Since $h^{m_i-1}\binom{n_i-1}{m_i-1}=\frac{t_i^{m_i-1}}{(m_i-1)!}
\prod_{1\le j<m_i}(1-j/n_i)$, the $m$-th term of Lemma~\ref{lem:runs} is
$\kappa_mv_m$ with
\[
\kappa_m=\prod_{i\in F}\Bigl[\prod_{1\le j<m_i}\Bigl(1-\frac j{n_i}\Bigr)\Bigr]
(1-hq_i)^{n_i-m_i}e^{hq_in_i}.
\]
Let $n_{\min}=\min_in_i$. If all $m_i\le n_i$, then
$(1-hq_i)^{n_i-m_i}\le e^{-hq_i(n_i-m_i)}$ gives $\kappa_m\le e^{hq_{\max}M}$.
Also $\log(1-x)\ge-x-x^2$ for $0\le x\le\frac12$ gives
$(1-hq_i)^{n_i-m_i}e^{hq_in_i}\ge e^{-n_ih^2q_i^2}$, and the product of these
over $i$ is at least $e^{-q_{\max}^2T^2/N}$. The factors $1-j/n_i$ lie in
$[0,1]$, so their product is at least
$1-\sum_im_i^2/(2n_i)\ge1-M^2/(2n_{\min})$. If some $m_i>n_i$, then
$\kappa_m=0$, while $M\ge n_{\min}+1$ gives $M^2/(2n_{\min})\ge2$. In all
cases
\begin{equation}\label{eq:A-kappa}
\Bigl(1-\frac{M^2}{2n_{\min}}\Bigr)e^{-q_{\max}^2T^2/N}\le\kappa_m\le
e^{hq_{\max}M}.
\end{equation}

Let $w_m=v_m/\sum_{m'}v_{m'}$, a probability since $H(t)>0$ under (H). By
Lemma~\ref{lem:runs} the ratio in the lemma is
$\E_w\kappa_m$, so by~\eqref{eq:A-kappa} we need $\E_we^{hq_{\max}M}$ and
$\E_wM^2$. Since $W$ for $\sigma A_F$ is $\sigma^{M-1}W(m)$, the function
$\psi(c)=\log\E_we^{c(M-1)}=\log H_{e^c}(t)-\log H(t)$ is convex with
$\psi(0)=0$. So $\psi(c)\le c\,\psi(1)$ for $0\le c\le1$, and $\psi(1)\ge1$
since $M\ge2$. We show that $\psi(1)\le C_{\mathcal K}T$.
By~\eqref{eq:A-Hsigma} with $\sigma=e$ and $\theta=\theta(\alpha)$,
$H_e(T\alpha)\le e^{k-1}b(\theta')\prod_{i\ge2}\theta_i\,
e^{eT\langle\alpha,\theta\rangle}$. For the lower bound note that
$H(T\alpha)e^{-T\langle\alpha,\theta\rangle}=T^{-(k-1)}e^{TI_F(\alpha)}f_T(\alpha)$.
By Theorem~\ref{thm:local} this is at least
$\frac12T^{-(k-1)/2}C_F(\alpha)$ for $T\ge T_1$, and for $1\le T\le T_1$ it is
continuous and positive on $[1,T_1]\times\mathcal K$. So in both cases
$H(T\alpha)\ge c_{\mathcal K}T^{-(k-1)/2}e^{T\langle\alpha,\theta\rangle}$.
Since $\theta(\alpha)$ and $b(\theta'(\alpha))$ are bounded on $\mathcal K$,
this gives $\psi(1)\le C_{\mathcal K}T$.

Now let $x=hq_{\max}\le\frac12$. Then
$\E_we^{xM}=e^{x+\psi(x)}\le e^{x(1+C_{\mathcal K}T)}=1+O(T^2/N)$, because
$T^2/N\le1$. The Chernoff bound $\Prob_w(M-1\ge y)\le e^{\psi(1)-y}$ gives
\[
\E_w(M-1)^2=\int_0^\infty2y\,\Prob_w(M-1\ge y)\,dy
\le\psi(1)^2+\int_{\psi(1)}^\infty2ye^{\psi(1)-y}\,dy
=\psi(1)^2+2\psi(1)+2,
\]
so $\E_wM^2=O(T^2)$. Finally $n_{\min}\ge\varepsilon_{\mathcal K}N$ with
$\varepsilon_{\mathcal K}=\min_{\alpha\in\mathcal K}\min_i\alpha_i>0$. So
by~\eqref{eq:A-kappa}, $\E_w\kappa_m$ lies between
$e^{-q_{\max}^2T^2/N}(1-\E_wM^2/(2\varepsilon_{\mathcal K}N))$ and
$\E_we^{hq_{\max}M}$, and both are $1+O(T^2/N)$, uniformly for
$\alpha\in\mathcal K$ and $1\le T\le N^{1/2}$. \qed

\subsection{Proof of \texorpdfstring{Theorem~\ref{thm:first}(ii)}{the
first-order bounds} for reducible faces}\label{app:first}

Irreducible admissible faces satisfy (H) by
Lemma~\ref{lem:admissible}(b), and Theorem~\ref{thm:facemax} covers them. So
let $F$ be admissible and reducible, with components $C_1,\dots,C_r$. Let $a$
be the least positive rate $q_{ij}$, and take $N_0$ so large that $h\le
N^{-2/3}$ is as small as we
need. Note that $T^2/N\le1$. As in the proof of
Lemma~\ref{lem:admissible}, a path of positive probability with support $F$
visits each component in one block of consecutive times, and it passes from
one block to the next by one switch. We call the order of the blocks, the
first state and these switches the pattern of the path. There are finitely
many patterns.

\emph{Upper bound.} Fix a pattern whose $j$-th block lies in the component
$D_j$, starts at $y_j$ and ends at $x_j$, and let $\supp n=F$. Block $j$ has
length $N_j=\sum_{i\in D_j}n_i$. The paths with this pattern and composition
$n$ have total probability $\mu_{y_1}\prod_{j<r}hq_{x_jy_{j+1}}$ times the
product over $j$ of the weight of the block paths from $y_j$ to $x_j$ inside
$D_j$ with the right composition. This is at most the probability that
$X_1,\dots,X_{N_j}$ have that composition for the chain started at $y_j$. For
$C=D_j$ with $|C|\ge2$, Lemma~\ref{lem:upper} with start $y_j$ and
$\theta=q-\lambda_C\ones$ bounds it by
$\kappa_1h^{|C|-1}e^{-h\lambda_CN_j+\kappa_1h^2N_j}$, where $\kappa_1$ is the
largest of these finitely many constants. For $C=\{i\}$ it is
$(1-hq_i)^{N_j-1}\le e^{1/2-h\lambda_CN_j}$. Since
$\sum_j(|D_j|-1)+r-1=k_F-1$, $\sum_jN_j=N$ and
$\lambda_C\ge\lambda_F$ for every component, summing over the patterns gives
$P_N(n)\le\kappa_2(T/N)^{k_F-1}e^{-T\lambda_F+\kappa_1}$.

\emph{Lower bound.} By Lemma~\ref{lem:admissible}(a) we can number the
components so that there are $y_1\in C_1$ with $\mu_{y_1}>0$ and
$x_j\in C_j$, $y_{j+1}\in C_{j+1}$ with $q_{x_jy_{j+1}}>0$ for $j<r$. Take
any $x_r\in C_r$, choose $j^*$ with $\lambda_{C_{j^*}}=\lambda_F$ and
$|C_{j^*}|=k^*$, and put $m=\lceil N/T\rceil$, so that
$1\le hm\le2$.

\emph{Claim A.} There is $c_A>0$ such that for every component $C$ and all
$y,x\in C$, the chain started at $y$ stays in $C$ for its first $m$ visits,
visits every state of $C$ and has $X_m=x$ with probability at least $c_A$.
Take a walk $y=v_0,\dots,v_J=x$ through every state of $C$ with
$q_{v_sv_{s+1}}>0$ and $J\le|C|^2$. The paths that make exactly
these switches, at any $J$ of the $m-1$ transition times, have total
probability at least
$\binom{m-1}J(ha)^J(1-hq_{\max})^m\ge(1-Jh)^Ja^Je^{-4q_{\max}}/J!$, since
$1\le hm\le2$. Also $(1-Jh)^J\ge2^{-J}$ for small $h$.

\emph{Claim B.} There is $c_B>0$ such that for all $y,x\in C_{j^*}$ and
$\ell\ge2m$, the same event with $C_{j^*}$ and $\ell$ visits has probability
at least $c_Be^{-h\lambda_F\ell-h^2\lambda_F^2\ell}$. Let $w>0$
satisfy $Q_{C_{j^*}}w=-\lambda_Fw$ and let $B=I+hQ_{C_{j^*}}\ge0$. Then
$(B^s)_{zz'}$ is the probability of going from $z$ to $z'$ in $s$ steps
inside $C_{j^*}$, and
$\sum_{z'}(B^s)_{zz'}\ge(B^sw)_z/\max w\ge(1-h\lambda_F)^s\min w/\max w$. We
use Claim A for the first $m$ visits, from $y$ back to $y$, then
$s=\ell-2m+1$ steps inside $C_{j^*}$, then Claim A for the last $m$ visits,
ending at $x$. This gives at least
$c_A^2(\min w/\max w)(1-h\lambda_F)^\ell$. Also $\lambda_F\le q_{\max}$,
because the Perron root of $Q_{C_{j^*}}$ is at least each diagonal entry. So
$h\lambda_F\le\frac12$ and $\log(1-h\lambda_F)\ge-h\lambda_F-h^2\lambda_F^2$.

Give block $j\ne j^*$ the length $m$ and block $j^*$ the length
$N-(r-1)m$, which is at least $2m$ for $T\ge T_0=r+2$ and $N\ge N_0$. The
paths that go from $y_j$ to $x_j$ in block $j$ as in the claims and pass from
block $j$ to block $j+1$ by the switch $x_j\to y_{j+1}$ all have support $F$.
By the Markov property their total probability is at least
$\mu_{y_1}(ha)^{r-1}c_A^{r-1}c_Be^{-T\lambda_F-q_{\max}^2}$. Their
compositions are fixed by the restrictions to the blocks, so there are at
most $(N+1)^{k^*-1}\prod_{j\ne j^*}(m+1)^{|C_j|-1}$ of them, and
$m+1\le3N/T$. Since $\sum_{j\ne j^*}(|C_j|-1)=k_F-k^*-(r-1)$,
dividing the probability by this number gives
$M_F\ge c\,N^{-(k_F-1)}T^{k_F-k^*}e^{-T\lambda_F}$. \qed
